\RequirePackage{pdfmanagement}
\SetKeys[document/metadata]{pdfversion=1.7, pdfstandard=A-2u, lang=en-GB}
\input{preamble/pdfA-compliance/glyphtounicode}
\documentclass[11pt,a4paper]{report}
\newcommand*{\PaperTitle}{Five reactant complexes are necessary and sufficient for absolute concentration robustness with bistability}
\newcommand*{\PaperTitlePlain}{Five reactant complexes are necessary and sufficient for absolute concentration robustness with bistability}
\newcommand*{\PaperShortTitle}{The minimum is five reactant complexes}
\newcommand*{\PaperSubtitle}{chemistry-open-problems \textperiodcentered{} 31 October 2026 \textperiodcentered{} verified}
\newcommand*{\PaperCredits}{Solver: GPT-6-Astra (effort xhigh) \textperiodcentered{} Verifier: Claude Opus 5.5 (effort not recorded)}
\newcommand*{\PaperKeywords}{absolute concentration robustness, bistability, reaction networks, exact enumeration, toric steady states}
\input{chemop-preamble}
\addbibresource{references.bib}
\providecommand{\R}{\mathbb R}
\newcommand{\Rp}{\mathbb R_{>0}}
\providecommand{\rank}{\operatorname{rank}}
\providecommand{\im}{\operatorname{im}}
\providecommand{\diag}{\operatorname{diag}}
\providecommand{\sgn}{\operatorname{sgn}}
\begin{document}
\makepapertitle
\begin{paperabstract}
The exact minimum number of distinct reactant complexes for coexistence of absolute concentration
robustness (ACR) and two exponentially stable positive equilibria in an at-most-bimolecular mass-action
network is five. The lower bound combines the earlier dimension-two obstruction with a complete
four-reactant exclusion: an exact finite classification reduces the latter to nineteen scalar normal
forms, each having at most two nondegenerate equilibria with opposite determinant signs. The upper
bound is an explicit four-species, six-reaction network with integer rates and three rational equilibria
in one compatibility class, two of which are stable. All computational claims have exact reproducible
certificates. The argument and the certificates were rechecked by a fresh-context verifier; there has
been no human peer review.
\end{paperabstract}

\section*{Question, result and scope}
Kaihnsa, Nguyen and Shiu~\autocite{KNS2024}, Section~6, write:
\begin{openquote}
``Another interesting direction, therefore, is to investigate conditions for coexistence of ACR and
multistability, rather than multistationarity.''
\end{openquote}
They conjecture stricter size bounds for this coexistence. The earlier repository manuscript~\autocite{Prior2026}
establishes dimension at least three, and bounds the number of reactant complexes between four and six.
Its remaining four-reactant cases, called S3, S7 and S9, were excluded only numerically.
Our precise follow-up question is the exact minimum number of reactant complexes. We determine it.

Here every source and every product has total molecularity at most two. The zero complex is allowed
and counts as a complex, also as a reactant complex: Example~4.13 of~\autocite{KNS2024} counts
$0\to X$ among its four reactant complexes. Rate constants are strictly positive. ACR means that at least one positive equilibrium exists and a fixed species has the same
concentration at \emph{all} positive equilibria, including equilibria in different classes. Stability is
exponential stability relative to the stoichiometric compatibility class. Arbitrarily many reactions from
each source are allowed; we do not assume one reaction per source.

\begin{theorem}\label{thm:minimum}
Among mass-action systems with at-most-bimolecular reactant and product complexes, the minimum number
of distinct reactant complexes allowing ACR and two exponentially stable positive equilibria in one
stoichiometric compatibility class is exactly five.
\end{theorem}
The lower and upper bounds concern arbitrary numbers of species and reactions. The five-source
construction uses four species; no claim of simultaneous minimization of all network sizes is made.

\begin{theorem}\label{thm:main}
Let a mass-action system have stoichiometric rank three, four distinct reactant complexes, and
at-most-bimolecular sources and products. If it has ACR, each positive compatibility class contains at most
two nondegenerate positive equilibria. When there are two, their restricted Jacobian determinants have
opposite signs. In particular, at most one is exponentially stable.
\end{theorem}

\begin{corollary}\label{cor:five}
An at-most-bimolecular mass-action system with ACR and two exponentially stable positive equilibria in one
class has at least five distinct reactant complexes.
\end{corollary}
\begin{proof}
With $r$ reactant complexes a positive equilibrium gives a strictly positive vector in the kernel of
the aggregated coefficient matrix $N$, so $\rank N\le r-1$. At a nondegenerate equilibrium,
$\im N$ equals the stoichiometric subspace $S$: indeed $\im J\subseteq\im N\subseteq S$, while
$J|_S$ is invertible. For $r\le4$ this forces $\dim S\le3$. Dimension two is excluded by the
dimension-two theorem of~\autocite{Prior2026}, which holds for every number of species and reactant
complexes. Dimension one is excluded by Proposition~3.10 of~\autocite{KNS2024}; directly, on a class the
field is a fixed vector times a polynomial of degree at most two in the class coordinate, and two
simple zeros of such a polynomial have opposite derivative signs. Dimension zero has singleton classes.
Dimension three forces $r=4$ and is excluded by \zcref{thm:main}.
\end{proof}
The rank-three theorem is proved here without using the earlier dimension-two theorem. Only the
corollary uses that earlier result. The following construction attains the bound and, together with
the corollary, proves \zcref{thm:minimum}.

\section{A five-source network attaining the minimum}
Consider four species $A,B,C,R$ and the six reactions
\begin{equation}\label{eq:network}
\begin{gathered}
 0\xrightarrow{64}C,\qquad A+B\xrightarrow{3}2B,\qquad
 A+C\xrightarrow{5}B+C,\qquad B+C\xrightarrow{64}A,\\
 0\xrightarrow{1}R,\qquad R\xrightarrow{1}0.
\end{gathered}
\end{equation}
Its five distinct reactant complexes are $0,A+B,A+C,B+C,R$. Every reactant and product has
molecularity at most two. With the mass-action convention $k x^y$ its equations are
\begin{equation}\label{eq:exampleode}
\begin{aligned}
 \dot a&=-3ab-5ac+64bc, &\quad \dot b&=3ab+5ac-64bc,\\
 \dot c&=64-64bc, &\quad \dot r&=1-r.
\end{aligned}
\end{equation}
The stoichiometric subspace is $\{v:v_a+v_b=0\}$, of dimension three; the only independent
conservation law is $a+b=T$. Every positive equilibrium has $r=1$, so ACR holds globally, across
all compatibility classes. In fact arbitrary positive rates in the last two reactions give the
ACR value $k_{0R}/k_{R0}$.

In the class $a+b=9$, equilibrium requires $b=9-a$, $c=1/(9-a)$ and $r=1$, with $0<a<9$.
The remaining equation factors exactly as
\[
 \dot a\big|_{b=9-a,\,c=1/(9-a)}
 =\frac{-3(a-4)(a-6)(a-8)}{9-a}=0.
\]
Thus the complete list of positive equilibria in this class is
\[
 P_1=(4,5,1/5,1),\qquad P_2=(6,3,1/3,1),\qquad P_3=(8,1,1,1).
\]
In the class coordinates $(a,c,r)$ the restricted Jacobian is
\[
 J_{\rm red}=\begin{pmatrix}
 6a-27-69c&576-69a&0\\
 64c&-64(9-a)&0\\
 0&0&-1
 \end{pmatrix}.
\]
Its characteristic polynomial is $(\lambda+1)(\lambda^2-\tau\lambda+\Delta)$, where
$\tau$ and $\Delta$ are the trace and determinant of the upper left $2\times2$ block:
\begin{center}
\begin{tabular}{crrl}
\toprule
Equilibrium & $\tau$ & $\Delta$ & Stability in the class\\
\midrule
$P_1$ & $-1684/5$ & $1536$ & Exponentially stable\\
$P_2$ & $-206$ & $-768$ & Saddle\\
$P_3$ & $-112$ & $1536$ & Exponentially stable\\
\bottomrule
\end{tabular}
\end{center}
The trace-negative, determinant-positive criterion for a real $2\times2$ matrix proves the two
stability assertions exactly; no numerical eigenvalue calculation is used.

The regulator can also affect the other species without increasing the number of sources. Replace
$0\xrightarrow{64}C$, $0\xrightarrow1 R$, $R\xrightarrow1 0$ by
$0\xrightarrow{62}C$, $0\xrightarrow1 C+R$, $R\xrightarrow1 C$.
Only $\dot c$ changes, by adding $r-1$. The equilibria are unchanged and the Jacobian remains
triangular in the two blocks, with exactly the same characteristic polynomial. Both versions are
verified directly from their reaction lists by \texttt{certify\_five\_example.py}.

\section{Toric reduction and the determinant sign}
List the sources as $y_1,\ldots,y_4$, and write $f(x)=N m(x)$, where $m_a(x)=x^{y_a}$.
It suffices to consider a system possessing a nondegenerate positive equilibrium $x^*$.
Then $\rank N=3$ and $\ker N$ is generated by $m(x^*)>0$.
Let $A$ be the $3\times n$ matrix with rows $y_a-y_4$, $a=1,2,3$.
The equilibrium equations are
\[
 A\log x=A\log x^*,\qquad
 \mathcal E=x^*\circ\exp(\ker A).
\]
Nondegeneracy forces $\rank A=3$. Otherwise the differential of these equations, and hence
the restricted Jacobian, has rank at most two. ACR in species $i$ is equivalent to
$e_i\in\operatorname{row}(A)$, or to a zero $i$th row in any basis matrix $V$ of $\ker A$.

Let $W$ have rows forming a basis of $S^\perp$. Compatibility is $Wx=Wx^*$.
Set $N_0=[N_1\ N_2\ N_3]$ and $p_a=m_a(x^*)/m_4(x^*)>0$.
Any three columns of $N$ are independent. At equilibrium,
\[
 J=m_4(x)N_0\diag(p_1,p_2,p_3)A\diag(1/x).
\]
Consequently, for a fixed nonzero constant $\sigma$ and a positive factor $q(x)$,
\begin{equation}\label{eq:orientation}
 \det(J|_S)=\sigma q(x)
 \det\begin{pmatrix}W\\ A\diag(1/x)\end{pmatrix}.
\end{equation}
To see this, use the basis $N_0$ for $S$ and complete it by columns $B$ to a basis of $\R^n$.
Multiplication of the displayed block matrix by $[B\ N_0]$ gives a block triangular matrix with
diagonal blocks $WB$ and $A\diag(1/x)N_0$. All change-of-basis determinants are fixed and nonzero.

The same determinant can be written with a basis matrix $V$ of $\ker A$. Choose $C$ with $AC$ invertible.
Then $[V\ C]$ is invertible, and multiplying the block matrix on the right by $\diag(x)[V\ C]$ gives a block
triangular matrix with diagonal blocks $W\diag(x)V$ and $AC$. Hence, for a fixed constant $c\ne0$,
\begin{equation}\label{eq:orientationV}
 \det\begin{pmatrix}W\\ A\diag(1/x)\end{pmatrix}
 =\frac{c}{x_1\cdots x_n}\,\det\bigl(W\diag(x)V\bigr).
\end{equation}
At a positive equilibrium, nondegeneracy is therefore equivalent to $\det(W\diag(x)V)\ne0$, and the sign
of $\det(J|_S)$ is a fixed sign times the sign of $\det(W\diag(x)V)$.

Replacing logarithmic binomials by ordinary binomials multiplies their differential rows by positive
monomial factors, followed by a fixed nonsingular row transformation. Fixing the ACR concentration and
performing linear eliminations therefore preserves the comparison of determinant signs.
In a one-variable reduction $F(t)=0$, the augmented determinant is a fixed sign times a positive
factor times $F'(t)$, provided the eliminated equations have nonsingular differentials on the positive
parameter interval. We verify this condition explicitly below. Two simple roots of a quadratic on an
interval have opposite derivative signs. The same holds for $h(t)-p$ when $p>0$ and $\log h$ is
strictly concave on that interval. Since a real Hurwitz $3\times3$ matrix has negative determinant,
opposite determinant signs preclude two stable equilibria.

\section{Exhaustive source classification}
A species absent from every source has a nonnegative derivative. At a positive equilibrium every
reaction rate is positive, so no reaction can produce that species either. Such species are inert and
can be removed. They cannot be ACR species, since their arbitrary positive concentrations give
equilibria with different values. Four bimolecular sources then involve at most eight species.

The source classification is elementary and finite. For $4\le n\le8$, enumerate four-element subsets
of
\[
 \{0,e_i,2e_i,e_i+e_j:1\le i<j\le n\}.
\]
Retain only sets covering all species, having $\rank A=3$, and having an ACR coordinate.
The rank test uses integer $3\times3$ minors. The ACR coordinates are precisely the indices contained
in every nonzero three-column minor: deleting such a column lowers the rank to two.
Canonicalization permutes the four sources and sorts the species columns, so no floating-point
rank decision or heuristic is involved. The script \texttt{certify\_sources.py} performs this enumeration.

For three species the toric equilibrium set is a singleton. For four species there are 26 source types.
In each, a generator $v$ of $\ker A$ has at most three distinct entries, including zero; this is checked
by its integer cofactor formula. The class equation is
\[
 \sum_{j=1}^4 W_j x_j^* e^{v_j t}-Wx^*=0.
\]
It has at most three distinct exponential terms. A nonzero sum of three real exponentials has at most
two real zeros counted with multiplicity: divide by the smallest exponential and apply Rolle's theorem
to the resulting derivative, which has at most two terms. If the sum is identically zero, all equilibria
on the curve are degenerate. In the nondegenerate case the two determinants have opposite signs by
\eqref{eq:orientation}.

For five or more species the exact list, up to species and source permutations, is as follows.
\begin{center}
\begin{tabular}{cll}
\toprule
Type & Sources & ACR species\\
\midrule
S1 & $0,E,C+D,A+B$ & $E$\\
S2 & $0,2E,C+D,A+B$ & $E$\\
S3 & $E,2E,C+D,A+B$ & $E$\\
S4 & $E,D,C+E,A+B$ & $C$\\
S5 & $E,D+E,C+E,A+B$ & $C,D$\\
S6 & $E,D+E,C+D,A+B$ & $D$\\
S7 & $E,D+E,2C,A+B$ & $D$\\
S8 & $E,D+E,B+C,A+C$ & $D$\\
S9 & $F,E+F,C+D,A+B$ & $E$\\
\bottomrule
\end{tabular}
\end{center}
There are eight types on five species, one on six species, and none on seven or eight.
This reproduces, in fully exact arithmetic, the earlier source list in~\autocite{Prior2026}.

For S4, S5, S6 and S8 the varying equilibrium concentrations, after positive constant factors are
removed, belong to three groups $g_1=u,g_2=w/u,g_3=w$. For S4, $c$ is fixed and $d,e$ are
proportional to $ab$; for S5, $c,d$ are fixed and $e$ is proportional to $ab$; for S6, $d$ is fixed
and $c,e$ are proportional to $ab$; for S8, $d$ is fixed, $a,b$ are proportional, and $e$ is
proportional to $ac$. The class equations are $Qg=T$, two linear equations in the three group
variables. If $\rank Q<2$, the equilibria are degenerate. Otherwise write the affine solution line
as $g(t)=g^0+t\rho$, where $\rho$ spans $\ker Q$. Positivity of all three components cuts out an
interval. The remaining equation is the polynomial $F(t)=g_1(t)g_2(t)-g_3(t)$ of degree at most two.
The augmented determinant with rows $Q$ and $\nabla(g_1g_2-g_3)$ is a fixed nonzero multiple of
$F'(t)$: the cross product of the rows of $Q$ is a fixed nonzero multiple of $\rho$.
The group embedding is linear with independent columns and positive constant entries on disjoint
coordinate groups. Eliminating the other fixed linear equations therefore changes the original
augmented determinant only by a fixed nonzero factor. Hence \eqref{eq:orientation} applies, and
two nondegenerate equilibria have opposite determinant signs. If $F$ vanishes identically, every
equilibrium is degenerate.

\section{Exhaustive stoichiometric classification}
We give a fresh exact classification for S1, S2, S3, S7 and S9. It supersedes the numerical subspace
search for the three previously unresolved cases.

For a fixed source set, form every possible nonzero reaction vector $p-y_a$, where $p$ is any
at-most-bimolecular product. Every rank-three stoichiometric subspace is spanned by three of these
vectors. Enumerate these triples and identify equal spans by their primitive exterior coordinates
(their $3\times3$ minors, divided by their gcd, with the first nonzero entry positive).
For each span $S$, retain all allowed reaction vectors lying in it, giving a maximal matrix $G$.
Membership is checked by the integer adjugate identity in a nonsingular three-coordinate projection.

\begin{lemma}\label{lem:maximal}
If some subnetwork with this span has a positive equilibrium, the maximal matrix $G$ has a strictly
positive kernel vector. In a three-coordinate projection, failure of this condition is detected by
a nonzero cross product of two columns having nonnegative inner products with every column,
after possibly reversing its sign.
\end{lemma}
\begin{proof}
The subnetwork equilibrium supplies a positive dependence among vectors spanning $S$. Their positive
cone is therefore all of $S$: any linear representation can be made nonnegative by adding a sufficiently
large multiple of that dependence. Each additional vector can thus be cancelled by a nonnegative
combination of the old vectors. Adding these dependencies to the original positive dependence gives
a dependence positive on every maximal-network reaction.

Conversely, a finite collection spanning $\R^3$ positively spans $\R^3$ exactly when its dual cone is
$\{0\}$. If this dual cone is nonzero it is pointed, since the original vectors span $\R^3$.
It has an extreme ray, on which at least two independent constraints are active. Such a ray is
perpendicular to two independent original vectors, and hence is their cross product up to sign.
This proves the exact finite test.
\end{proof}

Next write
\begin{equation}\label{eq:minors}
 D_W(z)=\det(W\diag(z)V)
       =\sum_{|I|=n-3}\det W_{:,I}\det V_{I,:}\prod_{i\in I}z_i.
\end{equation}
If all coefficients vanish, every equilibrium is degenerate by \eqref{eq:orientation} and
\eqref{eq:orientationV}. If the nonzero coefficients have one
sign, there is at most one positive equilibrium per class. Indeed, if $x,y$ are two such equilibria,
$\log x-\log y=Vt$ and $x-y=\diag(\ell)Vt$ with all logarithmic means $\ell_i>0$.
Then $W\diag(\ell)Vt=0$, contradicting \eqref{eq:minors} unless $t=0$.

A further exact exclusion concerns the aggregated kinetic rank. Let $K$ be a rational basis matrix
of $\ker G$, write a general flux as $K\theta$, and sum its reaction vectors separately for each
source, giving columns $L_1(\theta),\ldots,L_4(\theta)$ with sum zero. Project onto three independent
coordinates. If the cubic polynomial
\[
 \det[L_1(\theta)\ L_2(\theta)\ L_3(\theta)]
\]
vanishes identically, $\rank N<3$ at every positive equilibrium of every subnetwork with this span.
Such equilibria are degenerate. The certificate checks polynomial identity exactly. Nonzero evaluations
are used only to retain a case, never to discard one.

The complete counts are:
\begin{center}\small
\begin{tabular}{crrrrr}
\toprule
 & S1 & S2 & S3 & S7 & S9\\
\midrule
All rank-three spans & 4333&4287&4333&4397&23035\\
Missing a source &2965&3161&2965&3050&17734\\
No positive maximal flux &1126&782&1126&1042&4442\\
Injective &169&301&129&231&569\\
Class-degenerate &34&6&6&32&144\\
Kinetically degenerate &39&37&43&20&92\\
Remaining &0&0&64&22&54\\
\bottomrule
\end{tabular}
\end{center}
Each span occurs in exactly one row after the first row. After fixing the ACR coordinate, remove its
column from $W$ and row-reduce. Permutations $A\leftrightarrow B$, $C\leftrightarrow D$, and interchange
of the two pairs are allowed for S3 and S9; only $A\leftrightarrow B$ is needed for S7. The remaining
spaces reduce to eleven, five and three forms respectively. The next sections list \emph{all} of them
and prove the required root and determinant properties.

\section{The eleven S3 forms}
The positive equilibria have
\[
 e=e_0,\qquad ab=p,\qquad cd=q,\qquad p,q,e_0>0.
\]
After $e$ is fixed the two conservation equations have one of the following forms. Constants $T,U$
are arbitrary real numbers; existence of a positive solution supplies any needed positivity restrictions.
\begin{center}\small
\begin{tabular}{cll}
\toprule
Index & First equation & Second equation\\
\midrule
0 &$b=T$ &$c+d=U$\\
1 &$a-2c+2d=T$ &$b+c=U$\\
2 &$a-c+d=T$ &$b+(c+d)/2=U$\\
3 &$a-c+d=T$ &$b+c=U$\\
4 &$a-c+d=T$ &$b+c+d=U$\\
5 &$a-c+d=T$ &$b+2c=U$\\
6 &$a-c/2+d=T$ &$b+c/2=U$\\
7 &$a=T$ &$b+c+d=U$\\
8 &$a+d=T$ &$b+c=U$\\
9 &$a+(c+d)/2=T$ &$b+(c+d)/2=U$\\
10 &$a+c+d=T$ &$b+c+d=U$\\
\bottomrule
\end{tabular}
\end{center}
Forms 0 and 7 fix $a,b>0$ and $c+d$, so $cd=q$ is a quadratic on the interval $c,d>0$.
For form 8 take $a=t$, $b=p/t$, $c=U-p/t$, $d=T-t$. The positive interval is
$p/U<t<T$, with $T,U>0$. Multiplication of $cd-q=0$ by $t>0$ gives
\[
 -Ut^2+(UT+p-q)t-pT=0.
\]
For forms 9 and 10, subtraction fixes $a-b$; together with $ab=p$ this uniquely fixes the positive
pair $(a,b)$, since $a-p/a$ is strictly increasing from $-\infty$ to $+\infty$.
Again $c+d$ is fixed and $cd=q$ is quadratic on its positive interval.

For forms 1--6 put $c=t>0$, $d=q/t$. Then
\[
 a=T+\alpha t-\beta q/t,\qquad
 b=U-\gamma t-\delta q/t,
\]
where $\alpha,\beta,\gamma>0$ and $\delta\ge0$, as read directly from the table.
The functions $a,b$ are concave; $a$ is strictly concave. Their common positivity set is an interval.
On it,
\[
 (\log(ab))''=
 \frac{a''}{a}-\frac{(a')^2}{a^2}
 +\frac{b''}{b}-\frac{(b')^2}{b^2}<0.
\]
Hence $ab=p$ has at most two roots, with opposite derivative signs when two roots exist.

For completeness, take the augmented determinant with rows $R$, $\nabla(ab-p)$, $\nabla(cd-q)$,
where $R$ is the conservation matrix in the table. On the displayed scalar curves its multiplier of
the derivative is, up to a fixed sign: $T>0$ for forms 0 and 7; $t>0$ for forms 1--6 and 8;
and $h(a+b)>0$ for forms 9 and 10, where $h=1/2$ and $h=1$ respectively.
For form 8 multiplication of the scalar equation by $t$ preserves the sign at a root.
These identities also prove that repeated roots and identically zero reductions cannot be
nondegenerate equilibria.

\section{The five S7 forms}
Here the positive equilibrium equations are
\[
 d=d_0,\qquad ab=pe,\qquad c^2=qe,\qquad p,q,d_0>0.
\]
The reduced variables are $(a,b,c,e)$, and the conservation equations are:
\begin{center}\small
\begin{tabular}{cll}
\toprule
Index & First equation & Second equation\\
\midrule
0 &$b+e=T$ &$c-e=U$\\
1 &$a-c+2e=T$ &$b+c=U$\\
2 &$a-c/2+e/2=T$ &$b+c/2+e/2=U$\\
3 &$a-e=T$ &$b=U$\\
4 &$a+b+e=T$ &$c=U$\\
\bottomrule
\end{tabular}
\end{center}
For form 0 take $c=t$, $e=t-U$, $b=T-t+U$, $a=pe/b$.
Positivity is an interval, and $c^2-qe=t^2-q(t-U)=0$ is quadratic.
Its determinant multiplier has magnitude $b>0$.

For forms 1 and 2 put $u=1/c>0$, so $e=1/(qu^2)$. Then
\[
 A(u):=a/c=Tu+\alpha-\frac{\beta}{qu},\qquad
 B(u):=b/c=Uu-\gamma-\frac{\delta}{qu},
\]
with $(\alpha,\beta,\gamma,\delta)=(1,2,1,0)$ or $(1/2,1/2,1/2,1/2)$.
The positive interval is $A(u)>0$, $B(u)>0$. Both functions are concave and $A''(u)<0$.
Thus $\log(A B)$ is strictly concave, and $A B=p/q$ has at most two roots with opposite
derivative signs. If $H=AB-p/q$, the augmented determinant with rows $R$, $\nabla(ab-pe)$,
$\nabla(c^2-qe)$ equals, up to fixed orientation,
\[
 q\bigl(H'(u)-2H(u)/u\bigr).
\]
At a root its multiplier of $H'$ is the positive constant $q$.

For form 3 the equation $ab=pe$ is $UT+(U-p)e=0$, with $U=b>0$.
There is at most one positive solution, unless the equation is an identity, in which case the
positive solutions form a curve and are degenerate. For form 4, $c=U>0$ fixes $e=U^2/q$;
the sum and product of $a,b$ are fixed, giving a quadratic on $a,b>0$.
The determinant multiplier has magnitude $q>0$ in both scalar parametrizations used by the certificate.

\section{The three S9 forms}
The equilibrium equations are
\[
 e=e_0,\qquad ab=pf,\qquad cd=qf,\qquad p,q,e_0>0.
\]
In variables $(a,b,c,d,f)$ the three conservation equations have one of these forms:
\begin{center}\small
\begin{tabular}{clll}
\toprule
Index & First equation & Second equation & Third equation\\
\midrule
0 &$b+f=T$ &$c-f=U$ &$d=V$\\
1 &$a-c/2+f/2=T$ &$b+c/2+f/2=U$ &$d=V$\\
2 &$a=T$ &$b=U$ &$c+d+f=V$\\
\bottomrule
\end{tabular}
\end{center}
For form 0, put $f=t$, $b=T-t$, $c=t+U$, $d=V$, $a=pt/(T-t)$.
Positivity is an interval with $V>0$. The remaining equation is
$(V-q)t+VU=0$. It has at most one solution, or is an identity giving only degenerate equilibria.
The determinant multiplier has magnitude $b>0$.

For form 1, put $f=t$, $d=V>0$, $c=qt/V$,
\[
 a=T+\tfrac12(q/V-1)t,\qquad
 b=U-\tfrac12(q/V+1)t.
\]
The positive domain is an interval defined by three linear inequalities $t,a,b>0$.
The remaining equation $ab-pt=0$ is quadratic. Its determinant multiplier has magnitude $V>0$.

For form 2, $a=T>0$ and $b=U>0$ fix $f=TU/p$.
The sum $c+d=V-TU/p$ and product $cd=qTU/p$ are fixed, so there are at most two positive
solutions, with opposite derivative signs. The determinant multiplier has magnitude $p>0$.

All nineteen forms have now been treated. Together with the source classification, the injectivity
and degeneracy exclusions, and \eqref{eq:orientation}, this proves \zcref{thm:main}.

\section{Reproduction, checks and limitations}
From this directory, run:
\begin{verbatim}
python scripts/certify_sources.py
python scripts/certify_subspaces.py
python scripts/certify_normal_forms.py
python scripts/certify_five_example.py
\end{verbatim}
The first script re-enumerates every eligible source set. The second generates from scratch all
40,385 reaction spans for the five specified source types, proves each exclusion, and writes the
retained spaces and their classifications to \texttt{certificates/subspaces.json}. It does not read
the old numerical candidate lists. The third symbolically checks all nineteen determinant identities
and the concavity derivatives, including the conservation equations along every scalar curve.
The fourth derives the upper-bound example from its reaction list, factors its complete steady-state
equation and checks every stability claim in rational arithmetic.
Each script prints PASS/FAIL and exits nonzero on a failure.

The environment used was Python 3.12, SymPy 1.14 and NumPy under Fedora 44 in WSL. The runs take
less than one minute in total. Integer arrays use signed 64-bit entries; their input values are
in $[-2,2]$, three-minors have absolute value at most 48, adjugate entries at most 8, membership
products at most 288, and cross-product inner products at most 48. Thus these operations cannot
overflow. Nullspaces, polynomial identities and all subsequent coefficients are rational and symbolic.
No numerical eigenvalue, optimizer tolerance, random search or sampling result is used as an exclusion.

A fresh-context verifier reimplemented the enumeration, obtained the same label for every span, and recorded
two facts that make the proof sturdier; see \texttt{verification.md}. The positive-flux exclusion is
not needed: every span it rejects is also rejected by the injectivity or degeneracy tests. The kinetic
exclusion also has an elementary reason, without the cubic: on each span that it removes, the spans of the
reaction vectors of some three sources fail Rado's condition, so three columns of $N$ are forced to be dependent.

The primary source was read in its published HTML version. A public-source novelty audit on
31 October 2026, recorded in \texttt{references.md}, found no prior resolution of this exact-minimum question.
Absence of search hits is not a proof of priority. The manuscript is a computer-assisted proof, checked by
a fresh-context verifier, not by human referees. The five-reactant case is settled by \eqref{eq:network}.
The result does not concern larger products or nonhyperbolic notions of stability. It establishes
the exact reactant-complex minimum; it does not classify every network attaining it.
\paperreferences
\end{document}
